\section*{Chapter \ref{ch:g8:files-and-data-tables} --- Files and Data Tables}

\begin{solution}{exo:g8:files-and-data-tables:1}
\code{3 lines}, then the three lines, each without its \omterm{def:g8:files-and-data-tables:newline}{end of line}:
\code{lines[:4]} simply stops at the end of the \omterm{def:g7:lists:list}{list}.
\end{solution}

\begin{solution}{exo:g8:files-and-data-tables:2}
\code{['a', 'b', '', 'c']}: four fields, the third one empty. Two commas side
by side mean an empty field, like a missing temperature in the weather table.
\end{solution}

\begin{solution}{exo:g8:files-and-data-tables:3}
\code{str}: everything read from a \omterm{def:g8:files-and-data-tables:text-file}{text file} is text. \code{int(book["year"])}
gives the number.
\end{solution}

\begin{solution}{exo:g8:files-and-data-tables:4}
The counting pattern of \cref{ch:g8:dictionaries}, on one field of each
record:

\omcode{code/grade-8/03-files-and-data-tables/exercises/genres.py}{1}{6}{Books
per genre.}\label{lst:g8:files-and-data-tables:sol-genres}
\end{solution}

\begin{solution}{exo:g8:files-and-data-tables:5}
\code{book["pages"]} is text, and Python cannot add a number and a text. The
fix converts it: \code{total = total + int(book["pages"])}.
\end{solution}

\begin{solution}{exo:g8:files-and-data-tables:6}
Without \code{encoding}, Python~3.12 uses the encoding that the \omterm{def:g1:what-is-a-computer:computer}{computer} is
set up with, which differs from one \omterm{def:g1:what-is-a-computer:computer}{computer} to another. The \omterm{def:g3:files-and-folders:file}{file} is stored
in UTF-8, where \code{é} takes two \omterm{def:g4:memory-and-storage:byte}{bytes}; a \omterm{def:g1:what-is-a-computer:computer}{computer} set up with another
encoding reads those two \omterm{def:g4:memory-and-storage:byte}{bytes} as two \omterm{def:g7:strings:string}{characters}, \texttt{\~{A}\textcopyright}. Writing
\code{open("places.csv", encoding="utf-8")} reads it the same way everywhere.
\end{solution}

\begin{solution}{exo:g8:files-and-data-tables:7}
Count the loans per member, then keep the largest count; the \omterm{def:g7:functions:function}{function} returns
two \omterm{def:g6:variables-and-input:value}{values}:

\omcode{code/grade-8/03-files-and-data-tables/exercises/most_loans.py}{1}{10}{The
keenest reader.}\label{lst:g8:files-and-data-tables:sol-most-loans}
\end{solution}

\begin{solution}{exo:g8:files-and-data-tables:8}
\omcode{code/grade-8/03-files-and-data-tables/exercises/late.py}{1}{7}{Late
returns.}\label{lst:g8:files-and-data-tables:sol-late}

Text is compared \omterm{def:g7:strings:string}{character} by \omterm{def:g7:strings:string}{character} from the left, like words in a paper
\omterm{def:g8:dictionaries:dictionary}{dictionary}. In \code{YYYY-MM-DD} the year comes first, then the month, then
the day, each with a fixed number of digits: the first \omterm{def:g7:strings:string}{character} that
differs belongs to the largest unit that differs, exactly as when comparing
dates. Loans still out, with an empty return date, are left out.
\end{solution}

\begin{solution}{exo:g8:files-and-data-tables:9}
401 lines: the \omterm{def:g8:files-and-data-tables:table}{header row} and one line per record. With an \omterm{def:g8:files-and-data-tables:newline}{end of line} after
each, 401 end-of-line \omterm{def:g7:strings:string}{characters} too.
\end{solution}

\begin{solution}{exo:g8:files-and-data-tables:10}
\omcode{code/grade-8/03-files-and-data-tables/exercises/write_genre.py}{1}{11}{One
genre to a new \omterm{def:g3:files-and-folders:file}{file}.}\label{lst:g8:files-and-data-tables:sol-write}

\code{newline=""} lets the \code{csv} module write the ends of line itself.
A test of the chapter writes the poetry books and reads them back.
\end{solution}

\begin{solution}{exo:g8:files-and-data-tables:11}
\code{"09-02-2025"} (9 February) is smaller as text than
\code{"10-01-2025"} (10 January), because \code{0} comes before \code{1};
yet 9 February is later. With the day first, the first \omterm{def:g7:strings:string}{characters} compared
are the days, which matter least. With \code{YYYY-MM-DD} and a fixed number of
digits, the \omterm{def:g7:strings:string}{characters} are compared in the order of the units, from the
largest to the smallest.
\end{solution}

\begin{solution}{exo:g8:files-and-data-tables:12}
Keep the dates and the rain of the station, each date once; then try every
starting day and add 7 \omterm{def:g6:variables-and-input:value}{values} with a \omterm{def:g7:strings:slice}{slice}:

\omcode{code/grade-8/03-files-and-data-tables/exercises/rainiest_week.py}{1}{13}{The
rainiest 7 days.}\label{lst:g8:files-and-data-tables:sol-rainiest}

At Riverside it returns \code{('2025-01-13', 36.9)}: 36.9 millimetres from 13
to 19 January.
\end{solution}

\begin{solution}{pb:g8:files-and-data-tables:1}
\textbf{1.} On 1 January 2025, Hilltop had 3.9 degrees, no rain and a wind of
15.3 kilometres per hour; Riverside had 4.3 degrees, no rain and 10.4.

\textbf{2.} $2 \times 365 = 730$ rows, plus the \omterm{def:g8:files-and-data-tables:table}{header row}: 731 lines.

\textbf{3.} One line too many: a day read twice, for instance.

\textbf{4.} The months \code{"01"} to \code{"12"}, cut out of the date by the
\omterm{def:g7:strings:slice}{slice} \code{[5:7]}: the two \omterm{def:g7:strings:string}{characters} after \code{2025-}.

\textbf{5.} An empty field cannot be converted: \code{float("")} stops the
\omterm{def:g4:programs-and-operating-system:program}{program}. Python checks the left side of \code{and} first, and skips the right
side when the left is false.

\textbf{6.} March is the wettest, with 84.4 millimetres; 17 days are above 20
degrees:

\omoutput[1-4,14-15]{code/grade-8/03-files-and-data-tables/outputs/weather_year.txt}{The
wettest month and the warm days.}\label{out:g8:files-and-data-tables:weather-year}

\textbf{7.} Any month will do as a starting point, as long as it is a key of
\code{rain}; the walk then replaces it by any wetter month.

\textbf{8.} Missing temperatures, the error code $-99.0$, and dates seen
twice.

\textbf{9.}

\omoutput{code/grade-8/03-files-and-data-tables/outputs/defects.txt}{Riverside's
defects.}\label{out:g8:files-and-data-tables:defects}

366 rows for 365 days, because 14 June appears twice; 3 temperatures are
missing and 2 are the code $-99.0$.

\textbf{10.} An error code: the sensor failed, and the station wrote a number
that no real temperature can be, so that people notice.

\textbf{11.} From 12.04 to 11.43 degrees: about 0.6 degrees lower.

\textbf{12.} Each code is about 111 degrees below a normal \omterm{def:g6:variables-and-input:value}{value}. Spread over
363 \omterm{def:g6:variables-and-input:value}{values}, two such errors move the mean by about
$2 \times 111 / 363 \approx 0.6$ degrees.

\textbf{13.} \textbf{Result.} The wettest month at Hilltop in 2025 was March,
with 84.4 millimetres of rain. At Riverside, the two error codes of $-99.0$
pull the year's mean temperature from 12.04 down to 11.43 degrees, about 0.6
degrees.
\end{solution}
